\section*{Notation}
\addcontentsline{toc}{section}{Notation}

% Source: physical PDF pp. 21-23, printed pp. xx-xxii.

The big issue in choosing notation for a book on supersymmetry is whether to
use a two-component or a four-component notation for spinors. The standard
texts on supersymmetry have opted for the two-component Weyl notation. I have
chosen instead to use the four-component Dirac notation except in the first
stages of constructing the supersymmetry algebra and multiplets, because I
think this will make the book more accessible to those physicists who work on
particle phenomenology and model building. It would be a pity to see the
growth of a separate enclave of supersymmetry specialists, who communicate
well with each other but are cut off by their notation from the larger
community of particle theorists.

There is no great difficulty anyway in converting expressions in
four-component form into the two-component formalism. In the representation
of the Dirac matrices used throughout this book, in which \(\gamma_5\) is the
diagonal matrix with elements \(+1,+1,-1\), and \(-1\) on the main diagonal,
any four-component Majorana spinor \(\psi_\alpha\) (such as the supersymmetry
generator \(Q_\alpha\), the superspace coordinate \(\theta_\alpha\), or the
superderivative \(\mathcal{D}_\alpha\)) may be written in terms of a
two-component spinor \(\chi_a\) as
\[
  \psi=
  \begin{pmatrix}
    e\chi^*\\
    \chi
  \end{pmatrix},
\]
where \(e\) is the \(2\cdot2\) antisymmetric matrix with \(e_{12}=+1\). The
two-component spinor \(\chi_a\) is what in other books is often called
\(\psi_a^*=\bar{\psi}_{\dot a}\), while \((e\chi^*)_a\) would be called
\(\psi^a\). A summary of useful properties of four-component Majorana spinors
is given in the appendix to Chapter 26.

Here are some other features of the notation used in this book:

Latin indices \(i,j,k\), and so on generally run over the three spatial
coordinate labels, usually taken as \(1,2,3\). Where specifically indicated,
they run over values \(1,2,3,4\), with \(x^4=it\).

Greek indices \(\mu,\nu\), etc.\ from the middle of the Greek alphabet
generally run over the four spacetime coordinate labels \(0,1,2,3\), with
\(x^0\) the time coordinate. Where it is necessary to distinguish between
spacetime coordinates in a general coordinate system and in a locally inertial
system, indices \(\mu,\nu\), etc.\ are used for the former and \(a,b\), etc.\
for the latter.

Greek indices \(\alpha,\beta\), etc.\ from the beginning of the Greek alphabet
generally (except in Chapter 24) run over the components of four-component
spinors. To avoid confusion, I depart here from the conventions of Volume II,
and use upper-case letters \(A,B\), etc.\ to label the generators of a
symmetry algebra. Components of two-component spinors are labelled with
indices \(a,b\), etc. In particular, four-component supersymmetry generators
are denoted \(Q_\alpha\), while two-component generators (the bottom two
components of \(Q_\alpha\)) are called \(Q_a\).

Repeated indices are generally summed, unless otherwise indicated.

The spacetime metric \(\eta_{\mu\nu}\) is diagonal, with
\(\eta_{\mu\nu}=\operatorname{diag}(-1,+1,+1,+1)\).

The d'Alembertian is defined as
\(\Box\equiv\eta^{\mu\nu}\partial^2/\partial x^\mu\partial x^\nu
=\nabla^2-\partial^2/\partial t^2\), where \(\nabla^2\) is the Laplacian
\(\partial^2/\partial x^i\partial x^i\).

The `Levi--Civita tensor' \(\epsilon^{\mu\nu\rho\sigma}\) is defined as the
totally antisymmetric quantity with \(\epsilon^{0123}=+1\).

Dirac matrices \(\gamma^\mu\) are defined so that
\(\gamma^\mu\gamma^\nu+\gamma^\nu\gamma^\mu=2\eta^{\mu\nu}\). Also,
\(\gamma_5=i\gamma^0\gamma^1\gamma^2\gamma^3\), and
\(\beta=i\gamma^0=\gamma_4\). Where explicit matrices are needed, they are
given by the block matrices
\[
  \gamma^0=-i
  \begin{pmatrix}
    0 & \mathbf{1}\\
    \mathbf{1} & 0
  \end{pmatrix},
  \qquad
  \gamma^i=-i
  \begin{pmatrix}
    0 & \sigma^i\\
    -\sigma^i & 0
  \end{pmatrix},
\]
where \(\mathbf{1}\) is the unit \(2\cdot2\) matrix, \(0\) is the
\(2\cdot2\) matrix with elements zero, and the components of \(\boldsymbol
\sigma\) are the usual Pauli matrices
\[
  \sigma_1=
  \begin{pmatrix}
    0 & 1\\
    1 & 0
  \end{pmatrix},
  \qquad
  \sigma_2=
  \begin{pmatrix}
    0 & -i\\
    i & 0
  \end{pmatrix},
  \qquad
  \sigma_3=
  \begin{pmatrix}
    1 & 0\\
    0 & -1
  \end{pmatrix}.
\]
We also frequently make use of the \(4\cdot4\) block matrices
\[
  \gamma_5=
  \begin{pmatrix}
    \mathbf{1} & 0\\
    0 & -\mathbf{1}
  \end{pmatrix},
  \qquad
  \epsilon=
  \begin{pmatrix}
    e & 0\\
    0 & e
  \end{pmatrix},
\]
where \(e\) is again the antisymmetric \(2\cdot2\) matrix \(i\sigma_2\).
For instance, our phase convention for four-component Majorana spinors \(s\)
may be expressed as \(s^*=-\beta\gamma_5\epsilon s\).

The step function \(\theta(s)\) has the value \(+1\) for \(s>0\) and \(0\) for
\(s<0\).

The complex conjugate, transpose, and Hermitian adjoint of a matrix or vector
\(A\) are denoted \(A^*\), \(A^{\mathsf T}\), and
\(A^\dagger=(A^*)^{\mathsf T}\), respectively. We use a dagger for the
Hermitian adjoint of an operator and an asterisk for the complex conjugate of a
number. For a vector or matrix of operators, the dagger also denotes the
transpose of the matrix formed from the Hermitian adjoints of its components.
\(+\mathrm{H.c.}\) or \(+\mathrm{c.c.}\) at the end of an expression indicates
the addition of the Hermitian adjoint or complex conjugate of the foregoing
terms. A bar on a four-component spinor \(u\) is defined by
\(\bar u=u^\dagger\beta\).

Units are used with \(\hbar\) and the speed of light taken to be unity.
Throughout, \(-e\) is the rationalized charge of the electron, so that the
fine structure constant is \(\alpha=e^2/4\pi\simeq1/137\). Temperatures are in
energy units, with the Boltzmann constant taken equal to unity.

Numbers in parenthesis at the end of quoted numerical data give the
uncertainty in the last digits of the quoted figure. Where not otherwise
indicated, experimental data are taken from `Review of Particle Physics,' The
Particle Data Group, \emph{European Physics Journal C} \textbf{3}, 1 (1998).
